\frfilename{mt272.tex}
\versiondate{3.4.09}
\copyrightdate{2000}

\def\hatcalBnu{\hat\Cal B_{\nu}}

\def\chaptername{Teori probabilitas}
\def\sectionname{Kebebasan}

\newsection{272}

Saya memperkenalkan konsep 'kebebasan' bagi keluarga kejadian,
aljabar-sigma, dan variabel acak. Bagian pertama seksi ini, sampai 272G,
merupakan analisis hubungan-hubungan elementer di antara ketiga perwujudan
gagasan tersebut. Dalam 272G saya memberikan hasil mendasar bahwa distribusi
bersama suatu keluarga (berhingga) variabel acak yang bebas tidak lain
adalah produk distribusi masing-masing variabel. Ungkapan lebih lanjut tentang
hubungan antara kebebasan dan ukuran produk terdapat dalam 272J, 272M,
dan 272N. Saya memberikan suatu versi hukum nol-satu (272O), dan mengakhiri
seksi ini dengan sekumpulan hasil dasar dari teori probabilitas mengenai jumlah
dan hasil kali variabel-variabel acak yang bebas (272R-272W).
%272R 272S 272T 272U 272V 272W

\vleader{48pt}{272A}{Definisi} Misalkan $(\Omega,\Sigma,\mu)$ suatu
ruang probabilitas.

\header{272Aa}{\bf (a)} Suatu keluarga $\langle E_i\rangle_{i\in I}$ dalam
$\Sigma$ disebut {\bf bebas (secara stokastik)} jika

\Centerline{$\mu(E_{i_1}\cap E_{i_2}\cap\ldots\cap E_{i_n})
=\prod_{j=1}^n\mu E_{i_j}$}

\noindent setiap kali $i_1,\ldots,i_n$ adalah anggota-anggota berbeda dari $I$.

\header{272Ab}{\bf (b)} Suatu keluarga $\langle\Sigma_i\rangle_{i\in I}$ dari
subaljabar-sigma $\Sigma$ disebut {\bf bebas (secara stokastik)}
jika

\Centerline{$\mu(E_{1}\cap E_{2}\cap\ldots\cap E_{n})
=\prod_{j=1}^n\mu E_{j}$}

\noindent setiap kali $i_1,\ldots,i_n$ adalah anggota-anggota berbeda dari $I$
dan $E_j\in\Sigma_{i_j}$ untuk setiap $j\le n$.

\header{272Ac}{\bf (c)} Suatu keluarga $\langle X_i\rangle_{i\in I}$ dari
variabel acak bernilai real pada $\Omega$ disebut {\bf bebas (secara
stokastik)} jika

\Centerline{$\Pr(X_{i_j}\le\alpha_j\text{ untuk setiap }j\le n)
=\prod_{j=1}^n\Pr(X_{i_j}\le \alpha_j)$}

\noindent setiap kali $i_1,\ldots,i_n$ adalah anggota-anggota berbeda dari $I$
dan $\alpha_1,\ldots,\alpha_n\in\Bbb R$.

\cmmnt{
\leader{272B}{Catatan (a)} Barangkali inilah sumbangan utama teori
probabilitas kepada teori ukuran, dan karena itu patut dicermati dengan
sangat saksama. Gagasan 'kebebasan' berasal sepenuhnya dari luar
matematika, yakni dari gagasan tentang kejadian-kejadian yang mempunyai
sebab-sebab bebas. Saya kira 272G dan 272M merupakan hasil-hasil di
bawah ini yang paling jelas memperlihatkan aspek teori-ukur dari konsep
tersebut. Bukan suatu kebetulan bahwa keduanya melibatkan ukuran produk;
salah satu keajaiban teori ukuran ialah bahwa perangkat teknis yang sama
digunakan untuk membangun teori probabilitas tentang kebebasan
stokastik dan geometri volume multidimensi.

\header{272Bb}
{\bf (b)} Dalam paragraf-paragraf berikut saya akan mencoba menguraikan
beberapa hubungan antara ketiga pengertian kebebasan yang baru saja
didefinisikan. Namun patut segera diperhatikan bahwa, dalam ketiga kasus,
suatu keluarga bebas jika dan hanya jika semua subkeluarga
berhingganya bebas. Akibatnya, setiap subkeluarga dari suatu keluarga
bebas adalah bebas. Fakta elementer lain yang langsung mengikuti
definisi ialah bahwa jika $\langle \Sigma_i\rangle_{i\in I}$ merupakan
keluarga bebas dari aljabar-sigma, dan $\Sigma'_i$ merupakan
subaljabar-sigma dari $\Sigma_i$ untuk setiap $i$, maka
$\langle\Sigma'_i\rangle_{i\in I}$ merupakan keluarga bebas.

\header{272Bc}
{\bf (c)} Suatu perumusan ulang 272Ab yang berguna adalah sebagai berikut:
Suatu keluarga $\langle\Sigma_i\rangle_{i\in I}$ dari subaljabar-sigma
$\Sigma$ bebas jika dan hanya jika

\Centerline{$\mu(\bigcap_{i\in I}E_i)=\prod_{i\in I}\mu E_i$}

\noindent setiap kali $E_i\in\Sigma_i$ untuk setiap $i$ dan
$\{i:E_i\ne\Omega\}$ berhingga. (Di sini saya mengikuti konvensi 254F,
yakni bahwa untuk suatu keluarga $\langle\alpha_i\rangle_{i\in I}$ dalam
$[0,1]$ kita mengambil $\prod_{i\in I}\alpha_i=1$ jika $I=\emptyset$,
dan jika tidak, nilainya adalah
$\inf_{J\subseteq I,J\text{ berhingga}}\prod_{i\in J}\alpha_i$.)

\header{272Bd}
{\bf (d)} Dalam 272Aa-b saya berbicara tentang himpunan $E_i\in\Sigma$
dan aljabar $\Sigma_i\subseteq\Sigma$. Sebenarnya (272Ac sudah memberi
petunjuk mengenai hal ini), lebih sering daripada tidak kita akan
berurusan dengan $\hat\Sigma$, bukan $\Sigma$, jika keduanya berbeda,
dengan $(\Omega,\hat\Sigma,\hat\mu)$ pelengkapan dari
$(\Omega,\Sigma,\mu)$.
}%end of comment

\leader{272C}{Subaljabar-sigma yang didefinisikan oleh suatu variabel
acak}\cmmnt{ Untuk menghubungkan 272Ab dengan 272Ac kita memerlukan
pengertian berikut.} Misalkan $(\Omega,\Sigma,\mu)$ suatu ruang
probabilitas dan $X$ suatu variabel acak bernilai real yang didefinisikan
pada $\Omega$. Tuliskan $\Cal B$ untuk aljabar-sigma himpunan-himpunan
Borel dalam $\Bbb R$, dan $\Sigma_X$ untuk

\Centerline{$\{X^{-1}[F]:F\in\Cal B\}
\cup\{(\Omega\setminus\dom X)\cup X^{-1}[F]:F\in\Cal B\}$.}

\noindent Maka $\Sigma_X$ merupakan aljabar-sigma dari subhimpunan-subhimpunan
$\Omega$.
\prooflet{\Prf\

\Centerline{$\emptyset=X^{-1}[\emptyset]\in\Sigma_X$;}

\noindent jika $F\in\Cal B$, maka

\Centerline{$\Omega\setminus X^{-1}[F]
=(\Omega\setminus\dom X)\cup X^{-1}[\Bbb R\setminus F]\in\Sigma_X$,}

\Centerline{$\Omega\setminus((\Omega\setminus\dom X)\cup X^{-1}[F])
=X^{-1}[\Bbb R\setminus F]\in\Sigma_X$;}

\noindent jika $\sequence{k}{F_k}$ suatu barisan sebarang dalam $\Cal B$,
maka

\Centerline{$\bigcup_{k\in\Bbb N}X^{-1}[F_k]=X^{-1}[\bigcup_{k\in\Bbb
N}F_k]$,}

\noindent sehingga

\Centerline{$\bigcup_{k\in\Bbb N}X^{-1}[F_k]$,
\quad $(\Omega\setminus\dom X)\cup\bigcup_{k\in\Bbb N}X^{-1}[F_k]$}

\noindent termasuk dalam $\Sigma_X$. \Qed}

\cmmnt{Jelas bahwa }$\Sigma_X$ merupakan aljabar-sigma terkecil dari
subhimpunan-subhimpunan $\Omega$, yang memuat $\dom X$, dan yang terhadapnya
$X$ terukur. \cmmnt{Selain itu, }$\Sigma_X$ merupakan subaljabar dari
$\hat\Sigma$, dengan $\hat\Sigma$ domain pelengkapan $\mu$
\cmmnt{ (271Aa)}.

\cmmnt{Sekarang kita memperoleh hasil berikut.}

\leader{272D}{Proposisi} Misalkan $(\Omega,\Sigma,\mu)$ suatu ruang
probabilitas dan $\langle X_i\rangle_{i\in I}$ suatu keluarga variabel acak
bernilai real pada $\Omega$. Untuk setiap $i\in I$, misalkan $\Sigma_i$
aljabar-sigma yang didefinisikan oleh $X_i$\cmmnt{, seperti dalam 272C}.
Maka pernyataan-pernyataan berikut ekuivalen:

(i) $\langle X_i\rangle_{i\in I}$ bebas;

(ii) setiap kali $i_1,\ldots,i_n$ adalah anggota-anggota berbeda dari $I$
dan $F_1,\ldots,F_n$ merupakan subhimpunan Borel dari $\Bbb R$, maka

\Centerline{$\Pr(X_{i_j}\in F_j\text{ untuk setiap }j\le n)
=\prod_{j=1}^n\Pr(X_{i_j}\in F_j)$;}

(iii) setiap kali $\langle F_i\rangle_{i\in I}$ merupakan keluarga
subhimpunan Borel dari $\Bbb R$, dan $\{i:F_i\ne\Bbb R\}$ berhingga, maka

\Centerline{$\hat\mu\bigl(
\bigcap_{i\in I}(X_i^{-1}[F_i]\cup(\Omega\setminus\dom X_i))\bigr)
=\prod_{i\in I}\Pr(X_i\in F_i)$,}

\noindent dengan $\hat\mu$ pelengkapan $\mu$;

(iv) $\langle\Sigma_i\rangle_{i\in I}$ bebas terhadap $\hat\mu$.

\proof{{\bf (a)(i)$\Rightarrow$(ii)} Tuliskan
$\pmb{X}=(X_{i_1},\ldots,X_{i_n})$. Tuliskan $\nu_{\pmb{X}}$ untuk
distribusi bersama $\pmb{X}$, dan untuk setiap $j\le n$ tuliskan $\nu_j$
untuk distribusi $X_{i_j}$; misalkan $\nu$ produk
$\nu_1,\ldots,\nu_n$ sebagaimana diuraikan dalam 254A-254C. (Saya menulis
\S254 untuk produk tak berhingga. Jika Anda hanya berminat pada produk
berhingga ruang probabilitas, yang memadai untuk keperluan kita dalam
paragraf ini, saya menyarankan membaca \S\S251-252 dengan ketentuan dalam
pikiran bahwa semua ukuran adalah probabilitas, kemudian \S254 dengan
ketentuan bahwa himpunan $I$ berhingga.) Berdasarkan 256K, $\nu$ merupakan
ukuran Radon pada $\BbbR^n$. (Ini merupakan induksi pada $n$, dengan
mengandalkan 254N untuk memastikan bahwa kita dapat memandang $\nu$ sebagai
produk berulang
$(\ldots((\nu_1\times\nu_2)\times\nu_3)\times\ldots\nu_{n-1})
\times\nu_n$.) Maka untuk sebarang
$a=(\alpha_1,\ldots,\alpha_n)\in\BbbR^n$, kita mempunyai

$$\eqalignno{\nu\ocint{-\infty,a}
&=\nu\bigl(\prod_{j=1}^n\ocint{-\infty,\alpha_j}\bigr)
=\prod_{j=1}^n\nu_{j}\ocint{-\infty,\alpha_j}\cr
\noalign{\noindent (menggunakan 254Fb)}
&=\prod_{j=1}^n\Pr(X_{i_j}\le\alpha_j)
=\Pr(X_{i_j}\le\alpha_j\text{ untuk setiap }j\le n)\cr
\noalign{\noindent (menggunakan syarat (i))}
&=\nu_{\pmb{X}}\ocint{-\infty,a}.\cr}$$

\noindent Berdasarkan pernyataan ketunggalan dalam 271Ba,
$\nu=\nu_{\pmb{X}}$. Khususnya, jika $F_1,\ldots,F_n$ merupakan
subhimpunan Borel dari $\Bbb R$,

$$\eqalign{\Pr(X_{i_j}\in F_j\text{ untuk setiap }j\le n)
&=\Pr(\pmb{X}\in\prod_{j\le n}F_j)
=\nu_{\pmb{X}}(\prod_{j\le n}F_j)\cr
&=\nu(\prod_{j\le n}F_j)
=\prod_{j=1}^n\nu_{{j}}F_j
=\prod_{j=1}^n\Pr(X_{i_j}\in F_j),\cr}$$

\noindent sebagaimana diperlukan.


\medskip

{\bf (b)(ii)$\Rightarrow$(i)} langsung, jika kita mengingat bahwa semua
himpunan $\ocint{-\infty,\alpha}$ merupakan himpunan Borel, sehingga
definisi kebebasan yang diberikan dalam 272Ac hanyalah kasus khusus
dari (ii).

\medskip

{\bf (c)(ii)$\Rightarrow$(iv)} Andaikan (ii), dan misalkan
$i_1,\ldots,i_n$ adalah anggota-anggota berbeda dari $I$ dan
$E_j\in\Sigma_{{i_j}}$ untuk setiap $j\le n$. Untuk setiap $j$, tetapkan
$E'_j=E_j\cap\dom X_{i_j}$, sehingga $E'_j$ dapat dinyatakan sebagai
$X_{i_j}^{-1}[F_j]$ untuk suatu himpunan Borel $F_j\subseteq\Bbb R$.
Maka $\hat\mu(E_j\setminus E'_j)=0$ untuk setiap $j$, sehingga

$$\eqalignno{\hat\mu(\bigcap_{1\le j\le n}E_j)
&=\hat\mu(\bigcap_{1\le j\le n}E'_j)
=\Pr(X_{i_1}\in F_1,\ldots,X_{i_n}\in F_n)\cr
&=\prod_{j=1}^n\Pr(X_{i_j}\in F_j)\cr
\noalign{\noindent (menggunakan (ii))}
&=\prod_{j=1}^n\hat\mu E_j.\cr}$$

\noindent Karena $E_1,\ldots,E_n$ sebarang,
$\langle\Sigma_{i}\rangle_{i\in I}$ bebas.

\medskip

{\bf (d)(iv)$\Rightarrow$(ii)}
Sekarang andaikan $\langle\Sigma_{i}\rangle_{i\in I}$ bebas.
Jika $i_1,\ldots,i_n$ adalah anggota-anggota berbeda dari $I$ dan
$F_1,\ldots,F_n$ merupakan himpunan Borel dalam $\Bbb R$, maka
$X_{i_j}^{-1}[F_j]\in\Sigma_{i_j}$ untuk setiap $j$, sehingga

$$\eqalignno{\Pr(X_{i_1}\in F_1,\ldots,X_{i_n}\in F_n)
&=\hat\mu(\bigcap_{1\le j\le n}X_{i_j}^{-1}[F_j])\cr
&=\prod_{j=1}^n\hat\mu X_{i_j}^{-1}[F_j]
=\prod_{j=1}^n\Pr(X_{i_j}\in F_j)\cr.}$$

\medskip

%272D
{\bf (e)} Akhirnya, perhatikan bahwa (iii) hanyalah perumusan ulang dari
(ii), karena jika $F_i=\Bbb R$, maka $\Pr(X_i\in F_i)=1$ dan
$X_i^{-1}[F_i]\cup(\Omega\setminus\dom X_i)=\Omega$.
}%end of proof of 272D

\leader{272E}{Akibat} Misalkan $\langle X_i\rangle_{i\in I}$ suatu
keluarga bebas dari variabel acak bernilai real, dan
$\langle h_i\rangle_{i\in I}$ sebarang keluarga fungsi terukur Borel
dari $\Bbb R$ ke $\Bbb R$. Maka $\langle h_i(X_i)\rangle_{i\in I}$
bebas.

\proof{ Tuliskan $\Sigma_i$ untuk aljabar-sigma yang didefinisikan oleh
$X_i$, dan $\Sigma'_i$ untuk aljabar-sigma yang didefinisikan oleh
$h_i(X_i)$. Karena $h_i(X_i)$ terukur terhadap $\Sigma_i$ (121Eg), maka
$\Sigma'_i\subseteq\Sigma_i$ untuk setiap $i$ dan
$\langle\Sigma'_i\rangle_{i\in I}$ bebas, sebagaimana dalam 272Bb.
}%end of proof of 272E

\leader{272F}{}\cmmnt{ Dengan cara serupa, kita dapat menghubungkan
definisi dalam 272Aa dengan definisi-definisi lainnya.

\medskip

\noindent}{\bf Proposisi} Misalkan $(\Omega,\Sigma,\mu)$ suatu ruang
probabilitas, dan $\langle E_i\rangle_{i\in I}$ suatu keluarga dalam
$\Sigma$. Tetapkan
$\Sigma_i=\{\emptyset,E_i,\Omega\setminus E_i,\Omega\}$, yaitu
aljabar-sigma dari subhimpunan-subhimpunan $\Omega$ yang
dibangkitkan oleh $E_i$, dan $X_i=\chi E_i$, yaitu fungsi indikator
dari $E_i$. Maka pernyataan-pernyataan berikut ekuivalen:

(i) $\langle E_i\rangle_{i\in I}$ bebas;

(ii) $\langle \Sigma_i\rangle_{i\in I}$ bebas;

(iii) $\langle X_i\rangle_{i\in I}$ bebas.

\proof{{\bf (i)$\Rightarrow$(iii)} Jika $i_1,\ldots,i_n$ adalah
anggota-anggota berbeda dari $I$ dan
$\alpha_1,\ldots,\alpha_n\in\Bbb R$, maka untuk setiap $j\le n$ himpunan
$G_j=\{\omega:X_{i_j}(\omega)\le\alpha_j\}$ sama dengan salah satu dari
$E_{i_j}$, $\emptyset$, atau $\Omega$. Jika suatu $G_j$ kosong, maka

\Centerline{$\Pr(X_{i_j}\le\alpha_j\text{ untuk setiap }j\le
n\}=0=\prod_{j=1}^n\Pr(X_{i_j}\le\alpha_j)$.}

\noindent Jika tidak, tetapkan $K=\{j:G_j=E_{i_j}\}$; maka

$$\eqalign{\Pr(X_{i_j}\le\alpha_j\text{ untuk setiap }j\le n\}
&=\mu(\bigcap_{j\le n}G_j)
=\mu(\bigcap_{j\in K}E_{i_j})\cr
&=\prod_{j\in K}\mu E_{i_j}
=\prod_{j=1}^n\Pr(X_{i_j}\le\alpha_j).\cr}$$

\noindent Karena $i_1,\ldots,i_n$ dan $\alpha_1,\ldots,\alpha_n$
sebarang, $\langle X_i\rangle_{i\in I}$ bebas.

\medskip

{\bf (iii)$\Rightarrow$(ii)} mengikuti dari
(i)$\Rightarrow$(iii) dalam 272D, karena $\Sigma_i$ merupakan
aljabar-sigma yang didefinisikan oleh $X_i$.

\medskip

{\bf (ii)$\Rightarrow$(i)} langsung, karena $E_i\in\Sigma_i$ untuk setiap
$i$.
}%end of proof of 272F

\cmmnt{\medskip

\noindent{\bf Catatan} Saya berharap Anda akan merasa bahwa meskipun teori
ukuran produk mungkin sesuai untuk 272D, tentu teori tersebut merupakan
perangkat yang terlalu berat untuk digunakan pada sesuatu yang seharusnya
merupakan masalah kombinatorial sederhana seperti (iii)$\Rightarrow$(ii)
dalam proposisi ini. Saya menyarankan agar Anda menyusun bukti
'elementer', dan memeriksa gagasan-gagasan mana dari teori ukuran produk
(serta Teorema Kelas Monoton, 136B) yang sebenarnya diperlukan di sini.
}%end of comment

\leader{272G}{Distribusi variabel acak
bebas\dvrocolon{}}\cmmnt{ Saya belum mencoba menguraikan
'distribusi bersama' suatu keluarga tak berhingga dari variabel acak.
(Petunjuk tentang cara menangani keluarga terhitung diberikan dalam
271Ya dan 272Yb. Untuk keluarga tak terhitung saya akan menunggu sampai
\S454 dalam Jilid 4.) Namun, karena kebebasan suatu keluarga variabel
acak ditentukan oleh perilaku subkeluarga-subkeluarga berhingga, kita
dapat mendekatinya melalui proposisi berikut.

\medskip

\noindent}{\bf Teorema} Misalkan $\pmb{X}=(X_1,\ldots,X_n)$ suatu
keluarga berhingga variabel acak bernilai real pada suatu ruang
probabilitas. Misalkan $\nu_{\pmb{X}}$ distribusi yang bersesuaian pada
$\BbbR^n$. Maka pernyataan-pernyataan berikut ekuivalen:

(i) $X_1,\ldots,X_n$ bebas;

(ii) $\nu_{\pmb{X}}$ dapat dinyatakan sebagai produk dari $n$ ukuran
probabilitas $\nu_1,\ldots,\nu_n$, satu untuk setiap faktor $\Bbb R$
dari $\Bbb R^n$;

(iii) $\nu_{\pmb{X}}$ merupakan ukuran produk dari
$\nu_{X_1},\ldots,\nu_{X_n}$, dengan $\nu_{X_i}$ menyatakan distribusi
variabel acak $X_i$.

\proof{{\bf (a)(i)$\Rightarrow$(iii)} Dalam bukti
(i)$\Rightarrow$(ii) dari 272D di atas saya menunjukkan bahwa
$\nu_{\pmb{X}}$ merupakan produk $\nu$ dari
$\nu_{X_1},\ldots,\nu_{X_n}$.

\medskip

{\bf (b)(iii)$\Rightarrow$(ii)} langsung.

\medskip

{\bf (c)(ii)$\Rightarrow$(i)} Andaikan $\nu_{\pmb{X}}$ dapat
dinyatakan sebagai produk $\nu_1\times\ldots\times\nu_n$. Ambil
$a=(\alpha_1,\ldots,\alpha_n)$ dalam $\BbbR^n$. Maka

\Centerline{$\Pr(X_i\le\alpha_i\Forall i\le n)
=\Pr(\pmb{X}\in\ocint{-\infty,a})
=\nu_{\pmb{X}}(\ocint{-\infty,a})
=\prod_{i=1}^n\nu_i\ocint{-\infty,\alpha_i}$.}

\noindent Di sisi lain, dengan menetapkan
$F_i=\{(\xi_1,\ldots,\xi_n):\xi_i\le\alpha_i\}$, kita harus mempunyai

\Centerline{$\nu_i\ocint{-\infty,\alpha_i}
=\nu_{\pmb{X}}F_i=\Pr(\pmb{X}\in F_i)
=\Pr(X_i\le\alpha_i)$}

\noindent untuk setiap $i$. Jadi kita memperoleh

\Centerline{$\Pr(X_i\le\alpha_i\text{ untuk setiap }i\le n)
=\prod_{i=1}^n\Pr(X_i\le\alpha_i)$,}

\noindent sebagaimana diperlukan.
}%end of proof of 272G

\leader{272H}{Akibat} Andaikan $\langle X_i\rangle_{i\in I}$ suatu
keluarga bebas variabel acak bernilai real pada ruang probabilitas
$(\Omega,\Sigma,\mu)$, dan untuk setiap $i\in I$ diberikan suatu variabel
acak bernilai real lain $Y_i$ pada $\Omega$ sedemikian sehingga
$Y_i\eae X_i$. Maka $\langle Y_i\rangle_{i\in I}$ bebas.

\proof{ Untuk setiap $i_1,\ldots,i_n\in I$ yang berbeda, jika kita
menetapkan $\pmb{X}=(X_{i_1},\ldots,X_{i_n})$ dan
$\pmb{Y}=(Y_{i_1},\ldots,Y_{i_n})$, maka $\pmb{X}\eae\pmb{Y}$, sehingga
$\nu_{\pmb{X}}$, $\nu_{\pmb{Y}}$ sama (271De). Berdasarkan 272G,
$Y_{i_1},\ldots,Y_{i_n}$ harus bebas karena
$X_{i_1},\ldots,X_{i_n}$ bebas. Karena $i_1,\ldots,i_n$ sebarang,
seluruh keluarga $\langle Y_i\rangle_{i\in I}$ bebas.
}%end of proof of 272H

\cmmnt{\medskip

\noindent{\bf Catatan} Akibatnya kita dapat berbicara tentang keluarga
bebas dalam ruang $L^0(\mu)$ dari kelas-kelas ekuivalensi variabel
acak (241C), dengan mengatakan bahwa
$\langle X^{\ssbullet}_i\rangle_{i\in I}$ bebas jika dan hanya jika
$\langle X_i\rangle_{i\in I}$ bebas.
}%end of comment

\leader{272I}{Akibat} Andaikan $X_1,\ldots,X_n$ merupakan variabel acak
bernilai real yang bebas dengan fungsi-fungsi kepadatan
$f_1,\ldots,f_n$\cmmnt{ (271H)}.
Maka $\pmb{X}=(X_1,\ldots,X_n)$ mempunyai fungsi kepadatan $f$ yang
diberikan dengan menetapkan $f(x)=\prod_{i=1}^nf_i(\xi_i)$ setiap kali
$x=(\xi_1,\ldots,\xi_n)\in\prod_{i\le n}\dom(f_i)\subseteq\BbbR^n$.

\proof{ Untuk $n=2$ hal ini tercakup oleh 253I; kasus umum mengikuti
dengan induksi pada $n$.
}%end of proof of 272I

\leader{272J}{}\cmmnt{ Teorema-teorema terpenting dalam pokok bahasan ini
mengacu kepada keluarga bebas variabel acak, bukan keluarga
bebas aljabar-sigma. Nilai konsep aljabar-sigma bebas terletak
pada hasil-hasil seperti yang berikut.

\medskip

\noindent}{\bf Proposisi} Misalkan $(\Omega,\Sigma,\mu)$ suatu ruang
probabilitas lengkap, dan $\langle\Sigma_i\rangle_{i\in I}$ suatu
keluarga subaljabar-sigma dari $\Sigma$. Untuk setiap $i\in I$ misalkan
$\mu_i$ pembatasan $\mu$ pada $\Sigma_i$, dan misalkan
$(\Omega^I,\Lambda,\lambda)$ ruang probabilitas produk dari keluarga
$\langle(\Omega,\Sigma_i,\mu_i)\rangle_{i\in I}$. Definisikan
$\phi:\Omega\to\Omega^I$ dengan menetapkan $\phi(\omega)(i)=\omega$
setiap kali $\omega\in\Omega$ dan $i\in I$. Maka $\phi$ bersifat \imp\
jika dan hanya jika $\langle\Sigma_i\rangle_{i\in I}$ bebas.

\proof{ Ini hampir merupakan pernyataan ulang dari 254Fb dan 254G.
(i) Jika $\phi$ bersifat \imp, $i_1,\ldots,i_n\in I$ berbeda dan
$E_j\in\Sigma_{i_j}$ untuk setiap $j$, maka
$\bigcap_{j\le n}E_j=\phi^{-1}[\{x:x(i_j)\in E_j$ untuk setiap
$j\le n\}]$, sehingga

\Centerline{$\mu(\bigcap_{j\le n}E_j)
=\lambda\{x:x(i_j)\in E_j\text{ untuk setiap }j\le n\}
=\prod_{j=1}^n\mu_{i_j}E_j
=\prod_{j=1}^n\mu E_j$.}

\noindent (ii) Jika $\langle\Sigma_i\rangle_{i\in I}$ bebas,
$E_i\in\Sigma_i$ untuk setiap $i\in I$ dan
$\{i:E_i\ne\Omega\}$ berhingga, maka

\Centerline{$\mu\phi^{-1}[\prod_{i\in I}E_i]
=\mu(\bigcap_{i\in I}E_i)
=\prod_{i\in I}\mu E_i
=\prod_{i\in I}\mu_iE_i$.}

\noindent Jadi syarat-syarat 254G terpenuhi dan
$\mu\phi^{-1}[W]=\lambda W$ untuk setiap $W\in\Lambda$.
}%end of proof of 272J

\leader{272K}{Proposisi} Misalkan $(\Omega,\Sigma,\mu)$ suatu ruang
probabilitas dan $\langle\Sigma_i\rangle_{i\in I}$ suatu keluarga
bebas dari subaljabar-sigma $\Sigma$. Misalkan
$\langle J(s)\rangle_{s\in S}$ suatu keluarga saling lepas dari
subhimpunan-subhimpunan $I$, dan untuk setiap $s\in S$ misalkan
$\tilde\Sigma_s$ aljabar-sigma dari subhimpunan-subhimpunan $\Omega$
yang dibangkitkan oleh $\bigcup_{i\in J(s)}\Sigma_i$. Maka
$\langle\tilde\Sigma_s\rangle_{s\in S}$ bebas.

\proof{ Misalkan $(\Omega,\hat\Sigma,\hat\mu)$ pelengkapan dari
$(\Omega,\Sigma,\mu)$. Pada $\Omega^I$ misalkan $\lambda$ produk dari
ukuran-ukuran $\mu\restr\Sigma_i$, dan misalkan $\phi:\Omega\to\Omega^I$
peta diagonal, seperti dalam 272J. Berdasarkan 272J, $\phi$ bersifat
\imp\ untuk $\hat\mu$ dan $\lambda$.

Kita dapat mengidentifikasi $\lambda$ dengan produk
$\langle\lambda_s\rangle_{s\in S}$, dengan setiap
$\lambda_s\ (s\in S)$ merupakan produk dari
$\langle\mu\restr\Sigma_i\rangle_{i\in J(s)}$ (254N). Untuk $s\in S$,
misalkan $\Lambda_s$ domain $\lambda_s$, dan tetapkan
$\pi_s(x)=x\restr J(s)$ untuk $x\in\Omega^I$, sehingga $\pi_s$ bersifat
\imp\ untuk $\lambda$ dan $\lambda_s$ (254Oa), dan
$\phi_s=\pi_s\phi$ bersifat \imp\ untuk $\hat\mu$ dan $\lambda_s$;
tentu saja $\phi_s$ merupakan peta diagonal dari $\Omega$ ke
$\Omega^{J(s)}$. Tetapkan
$\Sigma_s^*=\{\phi_s^{-1}[H]:H\in\Lambda_s\}$. Maka $\Sigma_s^*$
merupakan subaljabar-sigma dari $\hat\Sigma$, dan
$\Sigma_s^*\supseteq\tilde\Sigma_s$, karena

\Centerline{$E=\phi_s^{-1}[\{x:x(i)\in E\}]\in\Sigma^*_s$}

\noindent setiap kali $i\in J(s)$ dan $E\in\Sigma_i$.

Sekarang andaikan $s_1,\ldots,s_n\in S$ berbeda dan
$E_j\in\tilde\Sigma_{s_j}$ untuk setiap $j$. Maka
$E_j\in\Sigma_{s_j}^*$, sehingga terdapat $H_j\in\Lambda_{s_j}$
sedemikian sehingga $E_j=\phi_{s_j}^{-1}[H_j]$ untuk setiap $j$.
Tetapkan

\Centerline{$W=\{x:x\in\Omega^I,\,x\restr J(s_j)\in H_j$ untuk setiap $j\le
n\}$.}

\noindent Karena kita dapat mengidentifikasi $\lambda$ dengan produk
dari $\lambda_s$, kita mempunyai

\Centerline{$\lambda W=\prod_{j=1}^n\lambda_{s_j}H_j
=\prod_{j=1}^n\hat\mu(\phi_{s_j}^{-1}[H_j])
=\prod_{j=1}^n\hat\mu E_j=\prod_{j=1}^n\mu E_j$.}

\noindent Di sisi lain, $\phi^{-1}[W]=\bigcap_{j\le n}E_j$, sehingga,
karena $\phi$ bersifat \imp,

\Centerline{$\mu(\bigcap_{j\le n}E_j)=\hat\mu(\bigcap_{j\le n}E_j)
=\lambda W=\prod_{j=1}^n\mu E_j$.}

\noindent Karena $E_1,\ldots,E_n$ sebarang,
$\langle\tilde\Sigma_s\rangle_{s\in S}$ bebas.
}%end of proof of 272K

\leader{272L}{}\cmmnt{ Saya memberikan suatu penerapan khas dari hasil
ini sebagai contoh.

\medskip

\noindent}{\bf Akibat} Misalkan $X,X_1,\ldots,X_n$ merupakan variabel
acak bernilai real yang bebas dan $h:\BbbR^n\to\Bbb R$ suatu fungsi
terukur Borel. Maka $X$ dan $h(X_1,\ldots,X_n)$ bebas.

\proof{ Misalkan $\Sigma_X$, $\Sigma_{X_i}$ aljabar-sigma yang
didefinisikan oleh $X$, $X_i$ (272C). Maka
$\Sigma_X,\Sigma_{X_1},\ldots,\Sigma_{X_n}$ bebas (272D).
Misalkan $\Sigma^*$ aljabar-sigma yang dibangkitkan oleh
$\Sigma_{X_1}\cup\ldots\cup\Sigma_{X_n}$. Maka 272K (barangkali dengan
bekerja dalam pelengkapan ruang probabilitas semula) memberi tahu kita
bahwa $\Sigma_X$ dan $\Sigma^*$ bebas. Namun setiap $X_j$ terukur
terhadap $\Sigma^*$, sehingga $Y=h(X_1,\ldots,X_n)$ terukur terhadap
$\Sigma^*$ (121Kb); selain itu $\dom Y\in\Sigma^*$, sehingga
$\Sigma_Y\subseteq\Sigma^*$ dan $\Sigma_X$, $\Sigma_Y$ bebas.
Berdasarkan 272D sekali lagi, $X$ dan $Y$ bebas, sebagaimana
dinyatakan.
}%end of proof of 272L

\cmmnt{\medskip

\noindent{\bf Catatan} Hampir semua dari kita, ketika mengajar teori
probabilitas elementer, akan mengajak mahasiswa memperlakukan akibat ini
(tentu saja dengan fungsi $h$ yang eksplisit) sebagai sesuatu yang
'jelas'. Pada hakikatnya, bukti di sini merupakan penegasan bahwa
definisi formal 'kebebasan' yang diberikan merupakan representasi
setia dari intuisi kita bahwa kejadian-kejadian bebas mempunyai
sebab-sebab bebas.
}%end of comment

\leader{272M}{}{\bf Produk ruang probabilitas dan keluarga bebas
variabel \dvrocolon{acak}}\cmmnt{ Kita telah melihat bahwa konsep
'variabel acak bebas' berhubungan erat dengan konsep 'ukuran
produk'. Sekarang saya memberikan beberapa perwujudan lebih lanjut dari
hubungan tersebut.

\medskip

\noindent}{\bf Proposisi} Misalkan
$\langle(\Omega_i,\Sigma_i,\mu_i)\rangle_{i\in I}$ suatu keluarga ruang
probabilitas, dan $(\Omega,\Sigma,\mu)$ produknya.

(a) Untuk setiap $i\in I$ tuliskan
$\tilde\Sigma_i=\{\pi_i^{-1}[E]:E\in\Sigma_i\}$, dengan
$\pi_i:\Omega\to\Omega_i$ peta koordinat. Maka
$\langle\tilde\Sigma_i\rangle_{i\in I}$ merupakan keluarga bebas
dari subaljabar-sigma $\Sigma$.

(b) Untuk setiap $i\in I$ misalkan
$\langle X_{ij}\rangle_{j\in J(i)}$ suatu keluarga bebas variabel
acak bernilai real pada $\Omega_i$, dan untuk $i\in I$, $j\in J(i)$
tuliskan $\tilde X_{ij}(\omega)=X_{ij}(\omega(i))$ bagi semua
$\omega\in\Omega$ sedemikian sehingga $\omega(i)\in\dom X_{ij}$. Maka
$\langle \tilde X_{ij}\rangle_{i\in I,j\in J(i)}$ merupakan keluarga
bebas variabel acak, dan setiap $\tilde X_{ij}$ mempunyai distribusi
yang sama dengan $X_{ij}$ yang bersesuaian.

\proof{{\bf (a)} Mudah diperiksa bahwa setiap $\tilde\Sigma_i$ merupakan
aljabar-sigma dari himpunan-himpunan. Selebihnya hanya mengingat kembali
dari 254Fb bahwa jika $J\subseteq I$ berhingga dan $E_i\in\Sigma_{i}$
untuk $i\in J$, maka

\Centerline{$\mu(\bigcap_{i\in J}\pi_i^{-1}[E_i])
=\mu\{\omega:\omega(i)\in E_i\text{ untuk setiap }i\in I\}
=\prod_{i\in I}\mu_{i}E_i$}

\noindent jika kita menetapkan $E_i=\Omega_i$ untuk $i\in I\setminus J$.

\medskip

{\bf (b)} Kita juga mengetahui bahwa $(\Omega,\Sigma,\mu)$ merupakan
produk dari pelengkapan-pelengkapan
$(\Omega_i,\hat\Sigma_i,\hat\mu_i)$ (254I).
Dari sini kita melihat bahwa setiap $\tilde X_{ij}$ terdefinisi
$\mu$-hampir di mana-mana, dan terukur terhadap $\Sigma$, dengan
distribusi yang sama dengan $X_{ij}$. Sekarang terapkan syarat (iii)
dari 272D. Andaikan
$\langle F_{ij}\rangle_{i\in I,j\in J(i)}$ suatu keluarga himpunan Borel
dalam $\Bbb R$, dan $\{(i,j):F_{ij}\ne\Bbb R\}$ berhingga. Tinjau

\Centerline{$E_i
=\bigcap_{j\in J(i)}
(X_{ij}^{-1}[F_{ij}]\cup(\Omega_i\setminus\dom X_{ij}))$,}

\Centerline{$E=\prod_{i\in I}E_i
=\bigcap_{i\in I,j\in J(i)}
(\tilde X_{ij}^{-1}[F_{ij}]\cup(\Omega\setminus\dom \tilde X_{ij}))$.}

\noindent Karena setiap keluarga $\langle X_{ij}\rangle_{j\in J(i)}$
bebas, dan $\{j:F_{ij}\ne\Bbb R\}$ berhingga,

\Centerline{$\hat\mu_iE_i=\prod_{j\in J(i)}\Pr(X_{ij}\in F_{ij})$}

\noindent untuk setiap $i\in I$. Karena

\Centerline{$\{i:E_i\ne \Omega_i\}\subseteq\{i:\exists\,j\in
J(i),\,F_{ij}\ne\Bbb R\}$}

\noindent berhingga,

\Centerline{$\mu E=\prod_{i\in I}\hat\mu_iE_i
=\prod_{i\in I,j\in J(i)}\Pr(\tilde X_{ij}\in F_{ij})$;}

\noindent karena $\langle F_{ij}\rangle_{i\in I,j\in J(i)}$ sebarang,
$\langle \tilde X_{ij}\rangle_{i\in I,j\in J(i)}$ bebas.
}%end of proof of 272M

\cmmnt{\medskip

\noindent{\bf Catatan} Perumusan dalam (b) lebih rumit daripada yang
diperlukan untuk menyatakan gagasannya, tetapi itulah yang diperlukan
untuk suatu penerapan di bawah.
}

\leader{272N}{}\cmmnt{ Suatu kasus khusus 272J sangat penting dalam teori
ukuran umum, dan paling berguna dalam bentuk yang disesuaikan.

\medskip

\noindent}{\bf Proposisi} Misalkan $(\Omega,\Sigma,\mu)$ suatu ruang
probabilitas lengkap, dan $\langle E_i\rangle_{i\in I}$ suatu keluarga
bebas dalam $\Sigma$ sedemikian sehingga $\mu E_i=\bover12$ untuk
setiap $i\in I$. Definisikan $\phi:\Omega\to\{0,1\}^I$ dengan menetapkan
$\phi(\omega)(i)=1$ jika $\omega\in E_i$, dan $0$ jika
$\omega\in\Omega\setminus E_i$. Maka $\phi$ merupakan fungsi pelestari
ukuran melalui prapeta untuk ukuran biasa $\lambda$ pada
$\{0,1\}^I$\cmmnt{ (254J)}.

\proof{ Saya menggunakan 254G sekali lagi. Untuk setiap $i\in I$
misalkan $\Sigma_i$ aljabar
$\{\emptyset,E_i,\Omega\setminus E_i,\Omega\}$; maka
$\langle\Sigma_i\rangle_{i\in I}$ bebas (272F). Untuk $i\in I$
tetapkan $\phi_i(\omega)=\phi(\omega)(i)$. Misalkan $\nu$ ukuran biasa
dari $\{0,1\}$. Maka mudah diperiksa bahwa

\Centerline{$\mu\phi_i^{-1}[H]=\Bover12\#(H)=\nu H$}

\noindent untuk setiap $H\subseteq\{0,1\}$.
Jika $\langle H_i\rangle_{i\in I}$ merupakan keluarga subhimpunan
$\{0,1\}$, dan $\{i:H_i\ne\{0,1\}\}$ berhingga, maka

$$\eqalignno{\mu\phi^{-1}[\prod_{i\in I}H_i]
&=\mu(\bigcap_{i\in I}\phi_i^{-1}[H_i])
=\prod_{i\in I}\mu\phi_i^{-1}[H_i]\cr
\noalign{\noindent (karena $\phi_i^{-1}[H_i]\in\Sigma_i$ untuk setiap $i$,
dan $\langle\Sigma_i\rangle_{i\in I}$ bebas)}
&=\prod_{i\in I}\nu H_i
=\lambda(\prod_{i\in I}H_i).\cr}$$

\noindent Karena $\langle H_i\rangle_{i\in I}$ sebarang, 254G memberikan
hasilnya.
}%end of proof of 272N

\leader{272O}{Aljabar-sigma ekor dan hukum
nol-satu\dvrocolon{}}\cmmnt{ Saya tidak pernah dapat memutuskan apakah
hasil berikut 'mendalam' atau tidak. Saya rasa ini merupakan salah satu
dari banyak kasus dalam matematika ketika suatu teorema mengejutkan dan
menggairahkan jika ditemukan tanpa persiapan, tetapi alamiah dan langsung
jika didekati dari sudut yang tepat.

\medskip

\noindent}{\bf Proposisi} Misalkan $(\Omega,\Sigma,\mu)$ suatu ruang
probabilitas dan $\sequencen{\Sigma_n}$ suatu barisan bebas dari
subaljabar-sigma $\Sigma$. Misalkan $\Sigma_n^*$ aljabar-sigma yang
dibangkitkan oleh $\bigcup_{m\ge n}\Sigma_m$ untuk setiap $n$, dan
tetapkan $\Sigma_{\infty}^*=\bigcap_{n\in\Bbb N}\Sigma_n^*$. Maka
$\mu E$ sama dengan $0$ atau $1$ untuk setiap
$E\in\Sigma_{\infty}^*$.

\proof{ Untuk setiap $n$, keluarga
$(\Sigma_0,\ldots,\Sigma_n,\Sigma^*_{n+1})$ bebas, berdasarkan 272K.
Jadi $(\Sigma_0,\ldots,\Sigma_n,\Sigma^*_{\infty})$ bebas, karena
$\Sigma^*_{\infty}\subseteq\Sigma^*_{n+1}$. Namun ini berarti bahwa
setiap subkeluarga berhingga dari
$(\Sigma^*_{\infty},\Sigma_0,\Sigma_1,\ldots)$ bebas, dan karena itu
seluruh keluarga tersebut juga bebas (272Bb). Akibatnya
$(\Sigma^*_{\infty},\Sigma^*_0)$ harus bebas, berdasarkan 272K
sekali lagi.

Sekarang jika $E\in\Sigma^*_{\infty}$, maka $E$ juga termasuk dalam
$\Sigma^*_0$, sehingga kita harus mempunyai

\Centerline{$\mu(E\cap E)=\mu E\cdot\mu E$,}

\noindent yaitu, $\mu E=(\mu E)^2$; jadi $\mu E\in\{0,1\}$, sebagaimana
dinyatakan.
}%end of proof of 272O

\leader{272P}{}\cmmnt{ Untuk mendukung klaim bahwa di suatu tempat kita
telah memperoleh wawasan yang tak trivial, saya memberikan suatu akibat,
yang akan menjadi dasar bagi pemahaman teorema-teorema limit dalam seksi
berikut, dan yang tampaknya tidak jelas.

\medskip

\noindent}{\bf Akibat} Misalkan $(\Omega,\Sigma,\mu)$ suatu ruang
probabilitas, dan $\langle X_n\rangle_{n\in\Bbb N}$ suatu barisan
bebas variabel acak bernilai real pada $\Omega$. Maka

\Centerline{$\limsup_{n\to\infty}\Bover{1}{n+1}(X_0+\ldots+X_n)$}

\noindent hampir di mana-mana konstan\dvro{.}{}\cmmnt{ -- yaitu, terdapat
suatu $u\in[-\infty,\infty]$ sedemikian sehingga

\Centerline{$\limsup_{n\to\infty}
\Bover{1}{n+1}(X_0+\ldots+X_n)=u$}

\noindent hampir di mana-mana.
}%end of comment

\proof{ Kita boleh mengandaikan bahwa setiap $X_n$ terukur terhadap
$\Sigma$ dan didefinisikan di setiap titik dalam $\Omega$, karena
(sebagaimana dicatat dalam 272H) mengubah $X_n$ pada suatu himpunan
terabaikan tidak memengaruhi kebebasannya, dan memengaruhi
$\limsup_{n\to\infty}\Bover1{n+1}(X_0+\ldots+X_n)$ hanya pada suatu
himpunan terabaikan. Untuk setiap $n$, misalkan $\Sigma_n$ aljabar-sigma
yang dibangkitkan oleh $X_n$ (272C), dan $\Sigma_n^*$ aljabar-sigma yang
dibangkitkan oleh $\bigcup_{m\ge n}\Sigma_m$; tetapkan
$\Sigma_{\infty}^*=\bigcap_{n\in\Bbb N}\Sigma_n^*$. Berdasarkan 272D,
$\sequencen{\Sigma_n}$ bebas, sehingga $\mu E\in\{0,1\}$ untuk
setiap $E\in\Sigma_{\infty}^*$ (272O).

Sekarang ambil sebarang $a\in\Bbb R$ dan tetapkan

\Centerline{$E_a=\{\omega:
\limsup_{m\to\infty}\Bover1{m+1}(X_0(\omega)+\ldots+X_m(\omega))
\le a\}$.}

\noindent Maka

\Centerline{$\limsup_{m\to\infty}\Bover1{m+1}(X_0+\ldots+X_m)
=\limsup_{m\to\infty}\Bover1{m+1}(X_n+\ldots+X_{m+n})$,}

\noindent sehingga


\Centerline{$E_a
=\{\omega:\limsup_{m\to\infty}
\Bover1{m+1}(X_n(\omega)+\ldots+X_{n+m}(\omega))\le a\}$}

\noindent termasuk dalam $\Sigma^*_n$ untuk setiap $n$, karena $X_i$
terukur terhadap $\Sigma_n^*$ untuk setiap $i\ge n$. Jadi
$E_a\in\Sigma_{\infty}^*$ dan

\Centerline{$\Pr(\limsup_{m\to\infty}
  \Bover1{m+1}(X_0+\ldots+X_m)\le a)
=\mu E_a$}

\noindent harus sama dengan $0$ atau $1$. Dengan menetapkan

\Centerline{$u=\sup\{a:a\in\Bbb R,\,\mu E_a=0\}$}

\noindent (dengan mengizinkan $\sup\emptyset=-\infty$ dan
$\sup\Bbb R=\infty$, seperti biasa dalam konteks semacam ini), kita
melihat bahwa

\Centerline{$\limsup_{n\to\infty}\Bover1{n+1}(X_0+\ldots+X_n)=u$}

\noindent hampir di mana-mana.
}%end of proof of 272P

\leader{*272Q}{}\cmmnt{ Saya menambahkan di sini suatu hasil yang akan
berguna dalam Jilid 5 dan yang memberikan wawasan lebih lanjut mengenai
sifat keluarga-keluarga bebas yang besar.

\medskip

\noindent}{\bf Teorema}\dvAnew{2009.}
Misalkan $(\Omega,\Sigma,\mu)$ suatu ruang probabilitas,
dan $\familyiI{\Sigma_i}$ suatu keluarga bebas dari subaljabar-sigma
$\Sigma$. Misalkan $\Cal E\subseteq\Sigma$ suatu keluarga himpunan
terukur, dan $\Tau$ aljabar-sigma yang dibangkitkan oleh $\Cal E$. Maka
terdapat suatu himpunan $J\subseteq I$ sedemikian sehingga
$\#(I\setminus J)\le\max(\omega,\#(\Cal E))$ dan
$\Tau$, $\family{j}{J}{\Sigma_j}$ bebas, dalam arti bahwa
$\mu(F\cap\bigcap_{r\le n}E_r)=\mu F\cdot\prod_{r=1}^n\mu E_r$ setiap
kali $F\in\Tau$, $j_1,\ldots,j_n$ adalah anggota-anggota berbeda dari
$J$ dan $E_r\in\Sigma_{j_r}$ untuk setiap $r\le n$.

\proof{{\bf (a)} Seperti dalam 272J, lengkapi $\Omega^I$ dengan ukuran
probabilitas $\lambda$ yang merupakan produk dari ukuran-ukuran
$\mu\restr\Sigma_i$, dan misalkan $\phi:\Omega\to \Omega^I$ peta
diagonal, sehingga $\phi$ bersifat \imp\ untuk $\hat\mu$ dan $\lambda$,
dengan $\hat\mu$ pelengkapan $\mu$.
Tuliskan $\Lambda$ untuk domain $\lambda$. Tetapkan
$\kappa=\max(\omega,\#(\Cal E))$, dan misalkan $\Cal E^*$ himpunan
$\{\bigcap_{r\le n}F_r:n\in\Bbb N$, $F_r\in\Cal E$ untuk setiap
$r\le n\}$. Karena $\#(\Cal E^n)\le\kappa$ untuk setiap $n$ (2A1Lc),
$\#(\Cal E^*)\le\kappa$ (2A1Ld). Untuk setiap
$F\in\Cal E^*$, definisikan $\nu_F:\Lambda\to[0,1]$ dengan menetapkan
$\nu_FW=\hat\mu(F\cap\phi^{-1}[W])$; maka $\nu_F$ aditif terhitung dan
didominasi oleh $\lambda$.
Karena itu fungsional tersebut mempunyai turunan Radon--Nikod\'ym $h_F$ terhadap
$\lambda$, sehingga $\hat\mu(F\cap\phi^{-1}[W])=\int_Wh_Fd\lambda$
untuk setiap $W\in\Lambda$ (232F). Berdasarkan 254Qc atau 254Rb, kita
dapat menemukan suatu fungsi $h'_F$ yang sama $\lambda$-hampir di
mana-mana dengan $h_F$ dan ditentukan oleh koordinat-koordinat
dalam suatu himpunan terhitung $J_F$, dalam arti bahwa
$h'_F(w)=h'_F(w')$ setiap kali $w$, $w'\in \Omega^I$ dan
$w\restr J_F=w'\restr J_F$. (Saya menerima begitu saja bahwa kita
memilih $h'_F$ agar didefinisikan di setiap titik pada $\Omega^I$.)

\medskip

{\bf (b)} Tetapkan $J=I\setminus\bigcup_{F\in\Cal E^*}J_F$; berdasarkan
2A1Ld, $I\setminus J=\bigcup_{F\in\Cal E^*}J_F$ mempunyai kardinalitas
paling besar $\kappa$.
Jika $F\in\Cal E^*$, $j_1,\ldots,j_n$ adalah anggota-anggota berbeda
dari $J$ dan $E_r\in\Sigma_{j_r}$ untuk setiap $r\le n$, maka
$\mu(F\cap\bigcap_{r\le n}E_r)=\mu F\cdot\prod_{r=1}^n\mu E_r$. \Prf\
Tetapkan $W=\{w:w\in \Omega^I$, $w(j_r)\in E_r$ untuk setiap
$r\le n\}$. Maka

\Centerline{$\mu(F\cap\bigcap_{r\le n}E_r)
=\hat\mu(F\cap\phi^{-1}[W])
=\int_Wh'_Fd\lambda
=\int h'_F\times\chi W\,d\lambda$.}

\noindent Namun perhatikan bahwa $W$ ditentukan oleh koordinat-koordinat
dalam $J$, sedangkan $h'_F$ ditentukan oleh koordinat-koordinat dalam
$J_F\subseteq I\setminus J$; dengan menggabungkan 272Ma, 272K, dan 272R
(atau dengan cara lain), kita mempunyai

\Centerline{$\mu(F\cap\bigcap_{r\le n}E_r)
=\int h'_F\times\chi W\,d\lambda
=\int h'_Fd\lambda\cdot\lambda W
=\mu F\cdot\prod_{r=1}^n\mu E_r$.  \Qed}

\medskip

{\bf (c)} Sekarang tinjau keluarga $\Cal A$ dari semua himpunan
$F\in\Sigma$ sedemikian sehingga
$\mu(F\cap\bigcap_{r\le n}E_r)
=\mu F\cdot\prod_{r=1}^n\mu E_r$ setiap kali
$j_1,\ldots,j_n\in J$ berbeda dan $E_r\in\Sigma_{j_r}$ untuk setiap
$r\le n$. Mudah diperiksa bahwa $\Cal A$ merupakan kelas Dynkin, dan
kita baru saja melihat bahwa $\Cal A$ memuat $\Cal E^*$; karena
$\Cal E^*$ tertutup terhadap $\cap$, $\Cal A$ memuat aljabar-sigma
$\Tau$ dari himpunan-himpunan yang dibangkitkan oleh $\Cal E^*$ (136B).
Dan inilah persis yang dinyatakan oleh teorema tersebut.
}%end of proof of 272Q

\leader{272R}{}\dvAformerly{2{}72Q}\cmmnt{ Sekarang saya harus mengejar
beberapa fakta dasar dari teori probabilitas elementer.

\medskip

\noindent}{\bf Proposisi} Misalkan $X$, $Y$ merupakan variabel acak
bernilai real yang bebas dengan ekspektasi
berhingga\cmmnt{ (271Ab)}. Maka $\Expn(X\times Y)$ ada dan sama dengan
$\Expn(X)\Expn(Y)$.

\proof{ Misalkan $\nu_{(X,Y)}$ distribusi bersama pasangan $(X,Y)$.
Maka $\nu_{(X,Y)}$ merupakan produk dari distribusi $\nu_X$ dan
$\nu_Y$ (272G). Selain itu $\int x\nu_X(dx)=\Expn(X)$ dan
$\int y\nu_Y(dy)=\Expn(Y)$ ada dalam $\Bbb R$ (271F). Jadi

\Centerline{$\int xy\nu_{(X,Y)}d(x,y)$ ada dengan nilai
$=\Expn(X)\Expn(Y)$}

\noindent (253D). Namun ini tidak lain adalah $\Expn(X\times Y)$,
berdasarkan 271E dengan $h(x,y)=xy$.
}%end of proof of 272R

\leader{272S}{Kesamaan Bienaym\'e}\dvAformerly{2{}72R} Misalkan
$X_1,\ldots,X_n$ merupakan variabel acak bernilai real yang bebas.
Maka
$\Var(X_1+\ldots+X_n)=\Var(X_1)+\ldots+\Var(X_n)$.

\proof{{\bf (a)} Pertama-tama andaikan semua $X_i$ mempunyai varians
berhingga. Tetapkan $a_i=\Expn(X_i)$, $Y_i=X_i-a_i$,
$X=X_1+\ldots+X_n$, dan $Y=Y_1+\ldots+Y_n$; maka
$\Expn(X)=a_1+\ldots+a_n$, sehingga $Y=X-\Expn(X)$ dan

$$\eqalignno{\Var(X)
&=\Expn(Y^2)
=\Expn(\sum_{i=1}^n Y_i)^2\cr
&=\Expn(\sum_{i=1}^n\sum_{j=1}^n Y_i\times Y_j)
=\sum_{i=1}^n\sum_{j=1}^n\Expn(Y_i\times Y_j).\cr}$$

\noindent Sekarang perhatikan bahwa jika $i\ne j$, maka
$\Expn(Y_i\times Y_j)=\Expn(Y_i)\Expn(Y_j)=0$, karena $Y_i$ dan $Y_j$
bebas (berdasarkan 272E) dan kita dapat menggunakan 272R, sedangkan
jika $i=j$, maka

\Centerline{$\Expn(Y_i\times Y_j)=\Expn(Y_i^2)=\Expn(X_i-\Expn(X_i))^2
=\Var(X_i)$.}

\noindent Jadi

\Centerline{$\Var(X)=\sum_{i=1}^n\Expn(Y_i^2)=\sum_{i=1}^n\Var(X_i)$.}

\medskip

{\bf (b)(i)} Selanjutnya saya menunjukkan bahwa jika
$\Var(X_1+X_2)<\infty$, maka $\Var(X_1)<\infty$. \Prf\ Kita mempunyai

$$\eqalignno{\iint(x+y)^2\nu_{X_1}(dx)\nu_{X_2}(dy)
&=\int (x+y)^2\nu_{(X_1,X_2)}(d(x,y))\cr
\noalign{\noindent (berdasarkan 272G dan teorema Fubini)}
&=\Expn((X_1+X_2)^2)\cr
\noalign{\noindent (berdasarkan 271E)}
&<\infty.\cr}$$

\noindent Jadi harus ada suatu $a\in\Bbb R$ sedemikian sehingga
$\int(x+a)^2\nu_{X_1}(dx)$ berhingga, yaitu,
$\Expn((X_1+a)^2)<\infty$; akibatnya $\Expn(X_1^2)$ dan $\Var(X_1)$
berhingga. \Qed

\medskip

\quad{\bf (ii)} Sekarang suatu induksi mudah (yang mengandalkan 272L!)
menunjukkan bahwa jika $\Var(X_1+\ldots+X_n)$ berhingga, maka
$\Var X_j$ juga berhingga untuk setiap $j$. Dengan membalikkan hal ini,
jika $\sum_{j=1}^n\Var(X_j)=\infty$, maka
$\Var(X_1+\ldots+X_n)=\infty$, dan sekali lagi keduanya sama.
}%end of proof of 272S

\leader{272T}{Distribusi jumlah variabel acak bebas:
Teorema}\dvAformerly{2{}72S} Misalkan
$X$, $Y$ merupakan variabel acak bernilai real yang bebas pada
ruang probabilitas $(\Omega,\Sigma,\mu)$, dengan distribusi $\nu_X$,
$\nu_Y$. Maka distribusi $X+Y$ adalah konvolusi
$\nu_X*\nu_Y$\cmmnt{ (257A)}.

\wheader{272T}{0}{0}{0}{48pt}

\proof{ Tetapkan $\nu=\nu_X*\nu_Y$. Ambil $a\in\Bbb R$ dan tetapkan
$h=\chi\ocint{-\infty,a}$. Maka $h$ terintegralkan terhadap $\nu$,
sehingga

$$\eqalignno{\nu\ocint{-\infty,a}
&=\int h\,d\nu
=\int h(x+y)(\nu_X\times\nu_Y)(d(x,y))\cr
\noalign{\noindent (berdasarkan 257B, dengan menuliskan
$\nu_X\times\nu_Y$ untuk ukuran produk pada $\BbbR^2$)}
&=\int h(x+y)\nu_{(X,Y)}(d(x,y))\cr
\noalign{\noindent (berdasarkan 272G, dengan menuliskan
$\nu_{(X,Y)}$ untuk distribusi bersama $(X,Y)$; di sinilah kita
menggunakan hipotesis bahwa $X$ dan $Y$ bebas)}
&=\Expn(h(X+Y))\cr
\noalign{\noindent (menerapkan 271E pada fungsi
$(x,y)\mapsto h(x+y)$)}
&=\Pr(X+Y\le a).\cr}$$

\noindent Karena $a$ sebarang, $\nu_X*\nu_Y$ merupakan distribusi
$X+Y$.
}%end of proof of 272T

\leader{272U}{Akibat}\dvAformerly{2{}72T} Andaikan $X$ dan $Y$ merupakan
variabel acak bernilai real yang bebas, dan keduanya mempunyai
kepadatan $f$ dan $g$. Maka konvolusi $f*g$ merupakan fungsi kepadatan
untuk $X+Y$.

\proof{ Berdasarkan 257F, $f*g$ merupakan fungsi kepadatan untuk
$\nu_X*\nu_Y=\nu_{X+Y}$.
}%end of proof of 272U

\leader{272V}{}\dvAformerly{2{}72U}\cmmnt{ Hasil sederhana berikut akan
sangat berguna ketika kita sampai pada proses stokastik dalam Jilid 4,
dan juga dalam seksi berikut.

\medskip

\noindent}{\bf Lemma Etemadi}\cmmnt{ ({\smc Etemadi 96})} Misalkan
$X_0,\ldots,X_n$ merupakan variabel acak bernilai real yang bebas.
Untuk $m\le n$, tetapkan $S_m=\sum_{i=0}^mX_i$. Maka

\Centerline{$\Pr(\sup_{m\le n}|S_m|\ge 3\gamma)
\le 3\max_{m\le n}\Pr(|S_m|\ge\gamma)$}

\noindent untuk setiap $\gamma>0$.

\proof{ Seperti dalam 272P, kita boleh mengandaikan bahwa setiap $X_i$
merupakan fungsi terukur yang didefinisikan di setiap titik pada suatu
ruang ukur $\Omega$. Tetapkan
$\alpha=\max_{m\le n}\Pr(|S_m|\ge\gamma)$. Untuk setiap $r\le n$,
tetapkan

\Centerline{$E_r=\{\omega:|S_m(\omega)|<3\gamma$ untuk setiap $m<r$,
$|S_r(\omega)|\ge 3\gamma\}$.}

\noindent Maka $E_0,\ldots,E_n$ merupakan partisi dari
$\{\omega:\max_{m\le n}|S_m(\omega)|\ge 3\gamma\}$. Tetapkan
$E'_r=\{\omega:\omega\in E_r$, $|S_n(\omega)|<\gamma\}$. Maka
$E'_r\subseteq\{\omega:\omega\in E_r$,
$|(S_n-S_r)(\omega)|>2\gamma\}$.
Namun $E_r$ bergantung pada $X_0,\ldots,X_r$, sehingga bebas dari
$\{\omega:|(S_n-S_r)(\omega)|>2\gamma\}$, yang dapat dihitung dari
$X_{r+1},\ldots,X_n$ (272K). Jadi

$$\eqalign{\mu E'_r
&\le\mu\{\omega:\omega\in E_r,\,|(S_n-S_r)(\omega)|>2\gamma\}
=\mu E_r\cdot\Pr(|S_n-S_r|>2\gamma)\cr
&\le\mu E_r(\Pr(|S_n|>\gamma)+\Pr(|S_r|>\gamma))
\le 2\alpha\mu E_r,\cr}$$

\noindent dan $\mu(E_r\setminus E'_r)\ge(1-2\alpha)\mu E_r$. Di sisi
lain, $\langle E_r\setminus E'_r\rangle_{r\le n}$ merupakan keluarga
saling lepas dari himpunan-himpunan yang semuanya termuat dalam
$\{\omega:|S_n(\omega)|\ge\gamma\}$. Jadi

\Centerline{$\alpha
\ge\mu\{\omega:|S_n(\omega)|\ge\gamma\}
\ge\sum_{r=0}^n\mu(E_r\setminus E'_r)
\ge(1-2\alpha)\sum_{r=0}^n\mu E_r$,}

\noindent dan

\Centerline{$\Pr(\sup_{r\le n}|S_r|\ge 3\gamma)
=\sum_{r=0}^n\mu E_r
\le\min(1,\Bover{\alpha}{1-2\alpha})
\le 3\alpha$,}

\noindent (dengan mempertimbangkan $\alpha\le\bover13$ dan
$\alpha\ge\bover13$ secara terpisah), sebagaimana diperlukan.
}%end of proof of 272V

\leader{*272W}{}\dvAnew{2009.}\cmmnt{ Hasil berikut merupakan penerapan
langsung yang serupa dari gagasan-gagasan seksi ini. Walaupun tidak akan
digunakan dalam jilid ini, hasil tersebut merupakan wakil yang mudah
dipahami dan berguna dari sangat banyak hasil mengenai ekor jumlah
variabel acak bebas.

\medskip

\noindent}{\bf Teorema}\cmmnt{ ({\smc Hoeffding 63})}
Misalkan $X_0,\ldots,X_n$ merupakan variabel acak bernilai real yang
bebas sedemikian sehingga $0\le X_i\le 1$ hampir di mana-mana\
untuk setiap $i$.
Tetapkan $S=\bover1{n+1}\sum_{i=0}^nX_i$ dan $a=\Expn(S)$. Maka

\Centerline{$\Pr(S-a\ge c)\le\exp(-2(n+1)c^2)$}

\noindent untuk setiap $c\ge 0$.

\proof{{\bf (a)} Tetapkan $a_i=\Expn(X_i)$ untuk setiap $i$.
Jika $b\ge 0$ dan $i\le n$, maka

\Centerline{$\Expn(e^{bX_i})\le\exp(ba_i+\Bover18b^2)$.}

\noindent\Prf\
Tetapkan $\phi(x)=e^{bx}$ untuk $x\in\Bbb R$. Maka $\phi$ konveks,
sehingga

\Centerline{$\phi(x)
\le 1+x(e^b-1)$}

\noindent setiap kali $x\in[0,1]$,

\Centerline{$\phi(X_i)
\leae 1+(e^b-1)X_i$}

\noindent dan

\Centerline{$\Expn(e^{bX_i})
=\Expn(\phi(X_i))
\le 1+(e^b-1)a_i
=e^{h(b)}$}

\noindent dengan $h(t)=\ln(1-a_i+a_ie^t)$
untuk $t\in\Bbb R$. Sekarang $h(0)=0$,

\Centerline{$h'(t)=\Bover{a_ie^t}{1-a_i+a_ie^t}
=1-\Bover{1-a_i}{1-a_i+a_ie^t}$,
\quad$h'(0)=a_i$,}

\Centerline{$h''(t)
=\Bover{1-a_i}{1-a_i+a_ie^t}\cdot\Bover{a_ie^t}{1-a_i+a_ie^t}
\le\Bover14$}

\noindent karena $a_ie^t$ dan $1-a_i$ keduanya lebih besar daripada atau
sama dengan $0$. Berdasarkan teorema Taylor dengan sisa, terdapat suatu
$t\in[0,b]$ sedemikian sehingga

\Centerline{$h(b)
=h(0)+bh'(0)+\Bover12b^2h''(t)
\le ba_i+\Bover18b^2$,}

\noindent dan

\Centerline{$\Expn(e^{bX_i})\le\exp(ba_i+\Bover18b^2)$.  \Qed}

\medskip

{\bf (b)} Ambil sebarang $b\ge 0$. Maka

$$\eqalignno{\Pr(S-a\ge c)
&=\Pr(\sum_{i=0}^n(X_i-a_i-c)\ge 0)
\le\Expn(\exp(b\sum_{i=0}^n(X_i-a_i-c)))\cr
\displaycause{karena $\exp(b\sum_{i=0}^n(X_i-a_i-c))\ge 1$ setiap kali
$\sum_{i=0}^n(X_i-a_i-c)\ge 0$}
&=e^{-(n+1)bc}\prod_{i=0}^ne^{-ba_i}\Expn(\prod_{i=0}^n\exp(bX_i))\cr
&=e^{-(n+1)bc}\prod_{i=0}^ne^{-ba_i}
  \prod_{i=0}^n\Expn(\exp(bX_i))\cr
\displaycause{karena variabel-variabel acak $\exp(bX_i)$ bebas,
berdasarkan 272E, sehingga ekspektasi hasil kali sama dengan hasil kali
ekspektasi-ekspektasinya, berdasarkan 272R}
&\le e^{-(n+1)bc}\prod_{i=0}^ne^{-ba_i}\exp(ba_i+\Bover18b^2)\cr
\displaycause{(a) di atas}
&=\exp(-(n+1)bc+\Bover{n+1}8b^2).\cr}$$

\noindent Sekarang nilai minimum dari fungsi kuadrat
$\Bover{n+1}8b^2-(n+1)cb$ adalah $-2(n+1)c^2$ ketika $b=4c$, sehingga
$\Pr(S-a\ge c)\le\exp(-2(n+1)c^2)$, sebagaimana dinyatakan.
}%end of proof of 272W

\exercises{
\leader{272X}{Latihan dasar (a)}
%\spheader 272Xa
Misalkan $(\Omega,\Sigma,\mu)$ suatu ruang probabilitas tanpa atom, dan
$\sequencen{\epsilon_n}$ sebarang barisan dalam $[0,1]$. Tunjukkan bahwa
terdapat suatu barisan bebas $\sequencen{E_n}$ dalam $\Sigma$
sedemikian sehingga $\mu E_n=\epsilon_n$ untuk setiap $n$.
\Hint{215D.}
%272B

\sqheader 272Xb Misalkan $\langle X_i\rangle_{i\in I}$ suatu keluarga
variabel acak bernilai real. Tunjukkan bahwa keluarga tersebut bebas
jika dan hanya jika

\Centerline{$\Expn(h_1(X_{i_1})\times\ldots\times h_n(X_{i_n}))
=\prod_{j=1}^n\Expn(h_j(X_{i_j}))$}

\noindent setiap kali $i_1,\ldots,i_n$ adalah anggota-anggota berbeda
dari $I$ dan $h_1,\ldots,h_n$ merupakan fungsi terukur Borel dari
$\Bbb R$ ke $\Bbb R$ sedemikian sehingga semua
$\Expn(h_j(X_{i_j}))$ berhingga.
%272E

\spheader 272Xc Tuliskan bukti 272F yang tidak menggunakan teori ukuran
produk.
%272F

\spheader 272Xd Misalkan $\pmb{X}=(X_1,\ldots,X_n)$ suatu keluarga
variabel acak bernilai real yang semuanya didefinisikan pada ruang
probabilitas yang sama, dan andaikan $\pmb{X}$ mempunyai fungsi kepadatan
$f$ yang dapat dinyatakan dalam bentuk
$f(\xi_1,\ldots,\xi_n)=f_1(\xi_1)f_2(\xi_2)\ldots f_n(\xi_n)$ untuk
fungsi-fungsi yang sesuai $f_1,\ldots,f_n$ dari satu variabel real.
Tunjukkan bahwa $X_1,\ldots,X_n$ bebas.
%272G

\spheader 272Xe Misalkan $X_1$, $X_2$ merupakan variabel acak bernilai
real yang bebas, keduanya dengan distribusi $\nu$ dan fungsi
distribusi $F$. Tetapkan $Y=\max(X_1,X_2)$. Tunjukkan bahwa distribusi
$Y$ kontinu mutlak terhadap $\nu$, dengan turunan Radon--Nikod\'ym
$F+F^{-}$, dengan $F^{-}(x)=\lim_{t\uparrow x}F(t)$ untuk setiap
$x\in\Bbb R$.
%272G

\spheader 272Xf Gunakan 254Sa dan gagasan 272J untuk memberikan bukti
lain bagi 272O.
%272O

\spheader 272Xg Misalkan $(\Omega,\Sigma,\mu)$ suatu ruang probabilitas
dan $\sequencen{\Sigma_n}$ suatu barisan tak menurun dari
subaljabar-sigma $\Sigma$. Misalkan $\Sigma_{\infty}$ aljabar-sigma yang
dibangkitkan oleh $\bigcup_{n\in\Bbb N}\Sigma_n$. Misalkan $\Tau$
subaljabar-sigma lain dari $\Sigma$ sedemikian sehingga $\Sigma_n$ dan
$\Tau$ bebas untuk setiap $n$. Tunjukkan bahwa $\Sigma_{\infty}$
dan $\Tau$ bebas. \Hint{terapkan Teorema Kelas Monoton pada
$\{E:\mu(E\cap F)=\mu E\cdot\mu F$ untuk setiap $F\in\Tau\}$.} Gunakan
hal ini untuk membuktikan 272O.
%272O

\spheader 272Xh Misalkan $\sequencen{X_n}$ suatu barisan variabel acak
bernilai real dan $Y$ suatu variabel acak bernilai real sedemikian
sehingga $Y$ dan $X_n$ bebas untuk setiap $n\in\Bbb N$. Andaikan
$\Pr(Y\in\Bbb N)=1$ dan
$\sum_{n=0}^{\infty}\Pr(Y\ge n)\Expn(|X_n|)$ berhingga. Tetapkan
$Z=\sum_{n=0}^YX_n$ (yaitu,
$Z(\omega)=\sum_{n=0}^{Y(\omega)}X_n(\omega)$ setiap kali
$\omega\in\dom Y$ sedemikian sehingga $Y(\omega)\in\Bbb N$ dan
$\omega\in\dom X_n$ untuk setiap $n\le Y(\omega)$). (i) Tunjukkan bahwa
$\Expn(Z)=\sum_{n=0}^{\infty}\Pr(Y\ge n)\Expn(X_n)$. \Hint{tetapkan
$X'_n(\omega)=X_n(\omega)$ jika $Y(\omega)\ge n$, dan $0$ jika tidak.}
(ii) Tunjukkan bahwa jika $\Expn(X_n)=\gamma$ untuk setiap $n\in\Bbb N$,
maka $\Expn(Z)=\gamma(\Expn(Y)+1)$. (Ini adalah {\bf persamaan Wald}
untuk konvensi jumlah yang dimulai dari indeks nol.)
%272R

\sqheader 272Xi Misalkan $X_1,\ldots,X_n$ merupakan variabel acak
bernilai real yang bebas. Tunjukkan bahwa jika
$X_1+\ldots+X_n$ mempunyai ekspektasi berhingga, maka demikian pula
setiap $X_j$. \Hint{bagian (b) dari bukti 272S.}
%272S

\sqheader 272Xj Misalkan $X$ dan $Y$ merupakan variabel acak bernilai
real yang bebas dengan kepadatan $f$ dan $g$. Tunjukkan bahwa
$X\times Y$ mempunyai fungsi kepadatan $h$, dengan
$h(x)=\int_{-\infty}^{\infty}\bover1{|y|}g(y)f(\bover{x}{y})dy$ untuk
hampir setiap $x$. \Hint{271K.}
%272U

\spheader 272Xk\dvAnew{2009.} Misalkan
$X_0,\ldots,X_n$ merupakan variabel acak bernilai real yang bebas
sedemikian sehingga $d_i\le X_i\le d'_i$ hampir di mana-mana\
untuk setiap $i$. (i) Tunjukkan bahwa jika $b\ge 0$, maka
$\Expn(e^{bX_i})\le\exp(ba_i+\Bover18b^2(d'_i-d_i)^2)$ untuk setiap $i$,
dengan $a_i=\Expn(X_i)$. (ii)
Tetapkan $S=\bover1{n+1}\sum_{i=0}^nX_i$ dan $a=\Expn(S)$. Tunjukkan
bahwa

\Centerline{$\Pr(S-a\ge c)\le\exp(-\Bover{2(n+1)^2c^2}{d})$}

\noindent untuk setiap $c\ge 0$, dengan
$d=\sum_{i=0}^n(d'_i-d_i)^2$.
%272W

\spheader 272Xl\dvAnew{2009.} Andaikan
$X_0,\ldots,X_n$ merupakan variabel acak bernilai real yang bebas,
semuanya dengan ekspektasi $0$, sedemikian sehingga
$\Pr(|X_i|\le 1)=1$ untuk setiap $i$. Tetapkan
$S=\Bover1{\sqrt{n+1}}\sum_{i=0}^nX_i$. Tunjukkan bahwa
$\Pr(S\ge c)\le\exp(-c^2/2)$ untuk setiap $c\ge 0$.
%272Xk 272W

\leader{272Y}{Latihan lanjutan (a)}
%\spheader 272Ya
\dvAformerly{2{}72Yc}
Misalkan $X_0,\ldots,X_n$ merupakan variabel acak bernilai real yang
bebas dengan distribusi $\nu_0,\ldots,\nu_n$ dan fungsi distribusi
$F_0,\ldots,F_n$. Tunjukkan bahwa, untuk sebarang himpunan Borel
$E\subseteq\Bbb R$,

\Centerline{$\Pr(\sup_{i\le n}X_i\in E)
=\sum_{i=0}^n\int_E\prod_{j=0}^{i-1}F_j^{-}(x)
  \prod_{j=i+1}^nF_j(x)\nu_i(dx)$,}

\noindent dengan $F_j^{-}(x)=\lim_{t\uparrow x}F_j(t)$ untuk setiap $j$,
dan kita menafsirkan hasil kali kosong
$\prod_{j=0}^{-1}F_j^{-}(x)$, $\prod_{j=n+1}^nF_j(x)$ sebagai $1$.
%272Xe, 272G

\spheader 272Yb\dvAformerly{2{}72Yf}
Misalkan $\pmb{X}=\sequencen{X_n}$ suatu barisan bebas variabel acak
bernilai real pada ruang probabilitas lengkap $(\Omega,\Sigma,\mu)$.
Misalkan $\Cal B$ aljabar-sigma Borel dari
$\BbbR^{\Bbb N}$ (271Ya). Misalkan $\nu_{\pmb{X}}^{(\Cal B)}$ ukuran
probabilitas dengan domain $\Cal B$ yang didefinisikan dengan menetapkan
$\nu_{\pmb{X}}^{(\Cal B)}E=\mu\pmb{X}^{-1}[E]$ untuk setiap
$E\in\Cal B$, dan tuliskan $\nu_{\pmb{X}}$ untuk pelengkapan dari
$\nu_{\pmb{X}}^{(\Cal B)}$. Tunjukkan bahwa $\nu_{\pmb{X}}$ tidak lain
adalah produk distribusi-distribusi $\nu_{X_n}$.
%272G

\spheader 272Yc\dvAformerly{2{}72Ye}
Misalkan $X_1,\ldots,X_n$ merupakan variabel acak bernilai real
sedemikian sehingga untuk setiap $j<n$ keluarga

\Centerline{$(X_1,\ldots,X_j,-X_{j+1},\ldots,-X_n)$}

\noindent mempunyai distribusi bersama yang sama dengan keluarga semula
$(X_1,\ldots,X_n)$. Tetapkan $S_j=X_1+\ldots+X_j$ untuk setiap
$j\le n$. (i) Tunjukkan bahwa untuk sebarang $a\ge 0$

\Centerline{$\Pr(\sup_{1\le j\le n}|S_j|\ge a)
\le 2\Pr(|S_n|\ge a)$.}

\noindent \Hint{tunjukkan bahwa jika
$E_j=\{\omega:\omega\in\bigcap_{i\le n}\dom X_i,\,|S_i(\omega)|<a$
untuk $i<j,\,|S_j(\omega)|\ge a\}$, maka
$\mu\{\omega:\omega\in E_j,\,
  |S_n(\omega)|\ge|S_j(\omega)|\}\ge\bover12\mu E_j$.}
(ii) Tunjukkan bahwa
$\Expn(\sup_{j\le n}|S_j|)\le 2\Expn(|S_n|)$. (iii)
Tunjukkan bahwa $\Expn(\sup_{i\le n}S_i^2)\le 2\Expn(S_n^2)$.
%272V

\spheader 272Yd\dvAnew{2009.}
Misalkan $\sequence{i}{X_i}$ suatu barisan bebas variabel acak
bernilai real, dan tetapkan $S_n=\sum_{i=0}^nX_i$ untuk setiap $n$.
Tunjukkan bahwa jika $\sequencen{S_n}$ konvergen ke $S$ dalam
$\eusm L^0$ untuk topologi konvergensi dalam ukuran, maka
$\sequencen{S_n}$ konvergen ke $S$ hampir di mana-mana.
%272V mt27bits

\spheader 272Ye\dvAformerly{2{}72Yd}
Misalkan $(\Omega,\Sigma,\mu)$ suatu ruang probabilitas.

\quad (i) Misalkan $\sequencen{E_n}$ suatu barisan bebas dalam
$\Sigma$. Tunjukkan bahwa untuk sebarang variabel acak bernilai real $X$
dengan ekspektasi berhingga,

\Centerline{$\lim_{n\to\infty}\int_{E_n}X\,d\mu-\mu E_n\Expn(X)=0$.}

\noindent \Hint{misalkan $\Tau_0$ subaljabar $\Sigma$ yang dibangkitkan
oleh $\{E_n:n\in\Bbb N\}$ dan $\Tau$ subaljabar-sigma $\Sigma$ yang
dibangkitkan oleh $\{E_n:n\in\Bbb N\}$. Mulailah dengan meninjau
$X=\chi E$ untuk $E\in\Tau_0$, kemudian $X=\chi E$ untuk $E\in\Tau$.
Berpindahlah dari $\eusm L^1(\mu\restrp\Tau)$ ke $\eusm L^1(\mu)$ dengan
menggunakan ekspektasi bersyarat.}

\quad (ii) Misalkan $\sequencen{X_n}$ suatu barisan bebas variabel
acak bernilai real yang terintegralkan secara seragam pada $\Omega$.
Tunjukkan bahwa untuk sebarang variabel acak bernilai real terbatas $Y$,

\Centerline{$\lim_{n\to\infty}\Expn(X_n\times Y)-\Expn(X_n)\Expn(Y)=0$.}

\quad(iii) Andaikan $1<p\le\infty$ dan tetapkan $q=p/(p-1)$ (dengan
mengambil $q=1$ jika $p=\infty$). Misalkan $\sequencen{X_n}$ suatu
barisan bebas variabel acak bernilai real dengan
$\sup_{n\in\Bbb N}\|X_n\|_p<\infty$, dan $Y$ suatu variabel acak bernilai
real dengan $\|Y\|_q<\infty$. Tunjukkan bahwa

\Centerline{$\lim_{n\to\infty}\Expn(X_n\times Y)-\Expn(X_n)\Expn(Y)=0$.}

\spheader 272Yf\dvAformerly{2{}72Yg}
Misalkan $(\Omega,\Sigma,\mu)$ suatu ruang probabilitas dan
$\sequencen{Z_n}$ suatu barisan variabel acak pada $\Omega$ sedemikian
sehingga $\Pr(Z_n\in\Bbb N)=1$ untuk setiap $n$, dan
$\Pr(Z_m=Z_n)=0$ untuk semua $m\ne n$. Misalkan $\sequencen{X_n}$ suatu
barisan variabel acak bernilai real pada $\Omega$, semuanya dengan
distribusi yang sama $\nu$, dan bebas satu sama lain serta dari
$Z_n$, dalam arti bahwa jika $\Sigma_n$ aljabar-sigma yang didefinisikan
oleh $X_n$, dan $\Tau_n$ aljabar-sigma yang didefinisikan oleh $Z_n$,
dan $\Tau$ aljabar-sigma yang dibangkitkan oleh
$\bigcup_{n\in\Bbb N}\Tau_n$, maka
$(\Tau,\Sigma_0,\Sigma_1,\ldots)$ bebas. Tetapkan
$Y_n(\omega)=X_{Z_n(\omega)}(\omega)$ setiap kali ini terdefinisi, yaitu,
$\omega\in\dom Z_n$, $Z_n(\omega)\in\Bbb N$, dan
$\omega\in\dom X_{Z_n(\omega)}$. Tunjukkan bahwa $\sequencen{Y_n}$
merupakan barisan bebas variabel acak dan bahwa setiap $Y_n$
mempunyai distribusi $\nu$.

\spheader 272Yg\dvAformerly{2{}72Yb}
Tunjukkan bahwa semua gagasan seksi ini berlaku dengan cara yang sama
bagi variabel acak bernilai kompleks, dengan penyesuaian yang sesuai
(yang harus dirancang).

\spheader 272Yh\dvAformerly{2{}72Ya} Kembangkan teori kebebasan untuk
variabel acak yang mengambil nilai dalam $\BbbR^r$, dengan menelusuri
sebanyak mungkin gagasan seksi ini.
}%end of exercises

\endnotes{
\Notesheader{272} Seksi ini panjang karena dua alasan: saya mencoba
memuat hasil-hasil dasar yang terkait dengan salah satu konsep matematika
yang paling subur, dan sulit mengetahui kapan harus berhenti; serta saya
mencoba melakukannya dalam bahasa yang sesuai bagi teori ukuran abstrak,
dengan menekankan berbagai pembedaan yang berada di pinggiran
gagasan-gagasan esensial. Walaupun saya bersedia bersikap luwes mengenai
pertanyaan apakah huruf $X$ seharusnya menyatakan suatu ruang atau suatu
fungsi, beberapa penerapan hasil-hasil ini yang terpenting bagi saya
berada dalam konteks yang menuntut kejelasan mutlak mengenai domain
fungsi-fungsi kita. Karena itu perlu dibentuk suatu pendapat mengenai
hal-hal seperti apa sebenarnya aljabar-sigma yang didefinisikan oleh suatu
variabel acak (272C).

Pokok 272Q ialah bahwa keluarga $\Cal E$ tidak harus berkaitan dengan
keluarga $\familyiI{\Sigma_i}$ dengan cara apa pun, kecuali, tentu saja,
bahwa kita berurusan dengan himpunan-himpunan terukur. Yang perlu kita
ketahui hanyalah bahwa $I$ cukup besar dibandingkan dengan $\Cal E$;
misalnya, bahwa $\Cal E$ terhitung dan $I$ tak terhitung. Keluarga
$\family{j}{J}{\Sigma_j}$ sekarang merupakan sejenis 'ekor' dari
$\familyiI{\Sigma_i}$, yang dengan aman bebas dari aljabar-sigma
'kepala' yang dibangkitkan oleh $\Cal E$.

Tentu saja saya harus menekankan sekali lagi bahwa bukti-bukti seperti
dalam 272R-272S harus dipandang sebagai penegasan bahwa kita mempunyai
model teori probabilitas yang sesuai, bukan sebagai alasan untuk
memercayai bahwa hasil-hasil tersebut sahih dalam konteks statistik.
Demikian pula, 272T-272U dapat didekati melalui berbagai intuisi mengenai
variabel acak diskret dan variabel acak dengan kepadatan kontinu, dan
walaupun hasil-hasil umum yang anggun itu menyenangkan, hasil-hasil
tersebut lebih penting bagi matematikawan murni daripada bagi ahli
statistik. Namun saya menemui hambatan aneh dalam bukti 272S, ketika
menunjukkan bahwa jika $X_1+\ldots+X_n$ mempunyai varians berhingga, maka
setiap $X_j$ juga demikian. Kita telah mempelajari teori ukuran yang cukup
untuk menangani hal ini dengan mudah, tetapi hubungannya dengan intuisi
probabilistik biasa, baik di sini maupun dalam 272Xi, tetap tidak jelas
bagi saya.

Ada empat gagasan dalam 272W yang patut disimpan untuk penggunaan nanti.
Yang pertama adalah taksiran

\Centerline{$\Expn(e^{bX_i})\le 1-a_i+e^ba_i$}

\noindent dalam bagian (a), suatu cara kasar tetapi efektif untuk
menggunakan hipotesis bahwa $X_i$ terbatas. Yang kedua adalah penggunaan
teorema Taylor untuk menunjukkan bahwa
$1-a_i+e^ba_i\le\exp(ba_i+\bover18b^2)$.
Yang ketiga adalah taksiran

\Centerline{$\Pr(Y\ge 0)\le\Expn(e^{bY})$ jika $b\ge 0$}

\noindent yang digunakan dalam bagian (b); dan yang keempat adalah 272R.
Setelah itu kita hanya perlu cukup bertekad untuk mencapai 272Xk. Namun
bahkan kasus khusus 272W pun mencolok sekaligus berguna.
}%end of notes

\discrpage
